% Generated draft: requires mathematical and source review.
\documentclass[11pt]{article}
\usepackage[margin=25mm]{geometry}
\usepackage{fontspec}
\setmainfont{Latin Modern Roman}
\usepackage{amsmath,amssymb,mathtools,unicode-math}
\setmathfont{Latin Modern Math}
\usepackage{enumitem,needspace,xurl,hyperref}
\setlength{\parindent}{0pt}
\setlength{\parskip}{5pt}
\setlength{\emergencystretch}{3em}
\newcommand{\N}{\mathbb{N}}
\newcommand{\Z}{\mathbb{Z}}
\newcommand{\Q}{\mathbb{Q}}
\newcommand{\R}{\mathbb{R}}
\newcommand{\C}{\mathbb{C}}
\newcommand{\F}{\mathbb{F}}
\newcommand{\A}{\mathbb{A}}
\newcommand{\PP}{\mathbb{P}}
\newcommand{\E}{\mathbb{E}}
\newcommand{\eps}{\varepsilon}
\newcommand{\dd}{\,\mathrm{d}}
\begin{document}
\section*{UnsolvedMath ID: 800577}
\textbf{Source problem code:} catalogue entry
\subsection{Definitions and mathematical statement}
Give each set \(X\subseteq\R\) its subspace topology. We use the following standard cover conventions, which are not explicitly specified in the supplied excerpt. A countable open cover of \(X\) is a countable family \(\mathcal U\) of relatively open subsets of \(X\) such that \(\bigcup\mathcal U=X\) and \(X\notin\mathcal U\). An \(\omega\)-cover is such a cover for which every finite subset of \(X\) is contained in some member of \(\mathcal U\). A \(\gamma\)-cover is an infinite such cover for which each point of \(X\) belongs to all but finitely many members. Write \(\Omega(X)\) and \(\Gamma(X)\) for these classes, respectively.

The property \(\mathsf S_1(\Gamma,\Gamma)\) means that
\[
\forall (\mathcal U_n)_{n\in\N}\in\Gamma(X)^{\N}\quad
\exists (U_n)_{n\in\N}\quad
\bigl[(\forall n\in\N)\ U_n\in\mathcal U_n
\ \text{ and }\ \{U_n:n\in\N\}\in\Gamma(X)\bigr].
\]
For the background comparison, \(\mathsf S_{\mathrm{fin}}(\Omega,\Omega)\) means that for every sequence \((\mathcal U_n)_{n\in\N}\) of members of \(\Omega(X)\), there are finite subfamilies \(\mathcal F_n\subseteq\mathcal U_n\) such that \(\bigcup_{n\in\N}\mathcal F_n\in\Omega(X)\).

For \(f,g\in\N^{\N}\), write \(f\leq^*g\) when \(f(n)\leq g(n)\) for all but finitely many \(n\). Define the bounding and dominating numbers by
\[
\mathfrak b=\min\bigl\{|B|:B\subseteq\N^{\N},\ (\forall g\in\N^{\N})(\exists f\in B)\ f\not\leq^*g\bigr\},
\]
\[
\mathfrak d=\min\bigl\{|D|:D\subseteq\N^{\N},\ (\forall g\in\N^{\N})(\exists f\in D)\ g\leq^*f\bigr\}.
\]
For infinite subsets \(A,B\subseteq\N\), write \(A\subseteq^*B\) if \(A\setminus B\) is finite. The tower number \(\mathfrak t\) is the least cardinal \(\kappa\) for which there exists a sequence \((T_\alpha)_{\alpha<\kappa}\) of infinite subsets of \(\N\), strictly decreasing modulo finite sets, with no infinite pseudointersection. Thus, for \(\alpha<\beta<\kappa\), one has \(T_\beta\subseteq^*T_\alpha\) and \(T_\alpha\not\subseteq^*T_\beta\), and there is no infinite \(A\subseteq\N\) with \(A\subseteq^*T_\alpha\) for every \(\alpha<\kappa\). The notation \(\operatorname{cf}(\kappa)\) denotes the least cardinality of an unbounded subset of the ordinal \(\kappa\). For a property \(P\) of sets of reals, put
\[
\operatorname{non}(P)=\min\{|Y|:Y\subseteq\R\text{ and }Y\text{ does not satisfy }P\}.
\]

\textbf{Problem.} Does ZFC prove the existence of a set \(X\subseteq\R\) such that
\[
|X|=\mathfrak b
\qquad\text{and}\qquad
X\models\mathsf S_1(\Gamma,\Gamma)?
\]
Here \(\mathfrak b\) is computed in the same set-theoretic universe as \(X\).

\textbf{Supplied historical context.} The source states that dichotomic arguments give, in ZFC, a set of reals satisfying \(\mathsf S_1(\Gamma,\Gamma)\) of cardinality \(\mathfrak t\), citing \texttt{[alpha\_i]}, and a set of reals satisfying \(\mathsf S_{\mathrm{fin}}(\Omega,\Omega)\) of cardinality \(\operatorname{cf}(\mathfrak d)\), citing \texttt{[SFH]}. It also states
\[
\operatorname{non}(\mathsf S_1(\Gamma,\Gamma))=\mathfrak b,
\qquad
\operatorname{non}(\mathsf S_{\mathrm{fin}}(\Omega,\Omega))=\mathfrak d,
\]
and that it is consistent that \(\mathfrak b>\mathfrak t\) and \(\mathfrak d>\operatorname{cf}(\mathfrak d)\). These comparisons motivate the source's description of the preceding existence results as unsatisfactory. They are retained as supplied historical assertions, not as independently verified results or a current literature assessment. The exact cover conventions and the bibliography entries behind \texttt{[alpha\_i]}, \texttt{[SFH]}, and the context label \texttt{[ideals]} require source review.
\subsection{Short English statement}
Does ZFC prove that some set of reals of cardinality exactly \(\mathfrak b\) satisfies \(\mathsf S_1(\Gamma,\Gamma)\)? This means that from each cover in any sequence of open \(\gamma\)-covers one can select one member so that the selected family is again a \(\gamma\)-cover. The supplied registry labels the problem open; that label has not been independently verified.
\subsection{Sources}
\begin{itemize}
\item Supplied local UnsolvedMath registry record, permanent ID \(800577\), exact source problem code \texttt{catalogue entry}; collection: AMR Open Problem Lists. The record attributes the problem to Tsaban, with proposed year 2006, and labels its status \texttt{open}. These are supplied registry assertions.
\item Tsaban, \emph{Selection Principles and Special Sets of Reals: Open Problems} (2006), Problem 3. Supplied source URL: \url{https://arxiv.org/abs/math/0609601}. The record reports extraction by \texttt{arxiv\_latex} and access on \texttt{2026-10-16}; this is supplied provenance, not an access or verification performed for this draft.
\item The supplied statement cites \texttt{[alpha\_i]} for the \(\mathfrak t\)-sized example and \texttt{[SFH]} for the \(\operatorname{cf}(\mathfrak d)\)-sized example. No corresponding bibliography entries are included in the supplied record, so these labels cannot be resolved here.
\item The supplied title is \emph{Selection Principles in Mathematics --- Problem 3 (\texttt{[ideals]})}; the background also gives \texttt{[ideals]} as its attribution/context label. No bibliographic expansion of this label is supplied. The recorded upstream revision is \texttt{d371ee917ac9647fbb06cd6fc529deec6f3dda82}.
\end{itemize}
\end{document}
